% 本文件是 SiSR 的历史档案，已从正文分离。见仓库 archive/README。
% 不要把它 include 进 SiSR.tex。在本目录下用 xelatex 单独编译。
\documentclass[cn,11pt,chinese]{elegantbook}

\title{历史笔记：Rocky 8}
\subtitle{已从 SiSR 正文分离的安装记录}
\author{冷轩}
\date{Archived}
\extrainfo{Centos 停更后的 Rocky 8 安装笔记，不再作为现行推荐。}

\begin{document}

\maketitle
\frontmatter
\tableofcontents
\mainmatter

\part{Rocky 8}
\chapter{系统安装}
Centos 是历年来服务器操作系统，但是停止更新了。而Rocky操作系统可以说是改了名之后的Centos，用法跟Centos一样。具体情况请自行搜索。

\begin{itemize}
\item 第一步：前往官网下载操作系统；

\item 第二步：创建U盘系统安装盘；
\begin{itemize}
\item Several applications can help you create a bootable USB. We recommend using Etcher(\url{https://www.balena.io/etcher/}). Download the application for your system (Windows, macOS or Linux), install and run.

The setup is intuitive and easy:
\begin{enumerate}
\item Select the CentOS 7 ISO image.
\item Insert the USB flash.
\item Find the USB and select it in the Select drive step.
\item Click Flash.
\end{enumerate}

\item 需要注意的是，我试过几款创建启动U盘的软件，会存在一些问题，导致系统安装不成功。让我几经折腾，最终发现是启动盘的问题。

\item 最终选用的启动U盘制作软件是：Fedora Media Writer。这款软件还有个好处是能把启动U盘重置为正常U盘，因为不能仅仅通过格式化来重置，否则不识别。
\end{itemize}

\end{itemize}



\section{yum}


\section{dnf}


\section{错误排查}
\subsection{查看日志文件}
\verb|sudo less /var/log/messages|

\end{document}
